
\begin{frame}{Two files import nothing --- and that emptiness is the feature}

\begin{columns}[T]
\begin{column}{0.46\textwidth}
\footnotesize
\key{\texttt{tdm\_split.py}} imports \texttt{dataclasses}. \\
\key{\texttt{lds\_geometry.py}} imports one name from \texttt{tdm\_split}. \\
No \texttt{Solution}, no \texttt{KernelWriter}, no \texttt{tP}, no rocisa.

\medskip
Functions needing more than three scalars take a \key{duck-typed context} --- so the
\emph{declaration} can live at the caller and the \emph{arithmetic} lives here.

\medskip
\good{Consequence:} \texttt{region\_row\_elems(32,0,False,2)} can be called in a bare
\texttt{pytest} and asserted \texttt{== 16}. The test it replaced had to match \emph{source
text}, and broke the moment the wrong expression was corrected.

\medskip
\bad{Do not add an import to either file.} Both docstrings state that obligation.
\end{column}

\begin{column}{0.52\textwidth}
\centering
\begin{tikzpicture}[node distance=3mm,font=\scriptsize]
\node[box,fill=white] (lra) {\texttt{Components/}\\\texttt{LraTileAssignment}};
\node[box,fill=white,right=6mm of lra] (lr) {\texttt{Components/}\\\texttt{LocalRead}};
\node[box,below=9mm of lra] (kwa) {\texttt{KernelWriter}\\\texttt{Assembly}};
\node[box,right=6mm of kwa] (lv) {\texttt{Lowering/}\\\texttt{leaves.py}};
\node[gbox,below=9mm of kwa,xshift=12mm,minimum width=44mm]
   (geo) {\texttt{lds\_geometry.py} \; + \; \texttt{tdm\_split.py}\\[1pt]
          \tsub{zero project imports}};
\draw[ar] (lra) -- (geo);
\draw[ar] (lr) -- (geo);
\draw[ar] (kwa) -- (geo);
\draw[ar] (lv) -- (geo);
\node[wbox,below=6mm of geo,text width=52mm] (com)
   {\bad{rejected:} \texttt{Tensile/Common/} --- the right home by layering,
    but its \texttt{\_\_init\_\_} imports rocisa};
\draw[ar,draw=warn,dashed] (geo) -- (com);
\end{tikzpicture}

\smallskip
\footnotesize
Four emitters, three layers, \key{one derivation}. \texttt{Components/} reaches \emph{up} into
\texttt{Lowering/} on purpose. The one remaining option --- duplicating the formula --- is what
caused the defect two slides on.
\end{column}
\end{columns}
\end{frame}
